\chapter{Los Logaritmos como Transformaciones que Preservan Estructura}

En los capítulos precedentes, la invariancia bajo cambios de coordenadas fue
identificada como la característica definitoria de los objetos geométricos y
tensoriales.
Un vector, una forma bilineal, un tensor --- ninguno queda definido por sus
componentes numéricas, sino por las relaciones estructurales que preserva bajo
transformación.
El presente capítulo constata que ese mismo principio opera fuera del ámbito
lineal.
El logaritmo es un isomorfismo entre estructuras algebraicas: convierte la
multiplicación en adición preservando íntegramente las relaciones que importan.
El invariante no es una coordenada, una longitud ni una componente tensorial ---
es la propia correspondencia algebraica.

% ------------------------------------------------------------------
%\section*{La Idea Central: una Aplicación que Preserva Estructura}
% ------------------------------------------------------------------

\medskip

En su nivel más elemental, el logaritmo responde a la pregunta:
\emph{¿qué exponente produce un número dado?}
Sin embargo, tanto histórica como matemáticamente,
el poder de los logaritmos radica en algo más profundo:
convierten la multiplicación en adición --- y lo hacen de forma exacta.

\medskip

En el Capítulo~I se definió un isomorfismo como una aplicación entre sistemas
algebraicos que preserva su estructura.
El logaritmo proporciona una realización precisa y global de esa idea.

\medskip

La aplicación
\[
\log:\mathbb{R}^+\to\mathbb{R}
\]
no es una mera biyección; es un isomorfismo entre dos estructuras de grupo:

\begin{enumerate}
    \item Un grupo multiplicativo $(\mathbb{R}^+,\cdot)$.
    \item Un grupo aditivo $(\mathbb{R},+)$.
    \item Un isomorfismo
    \[
    \log:(\mathbb{R}^+,\cdot)\longrightarrow(\mathbb{R},+),
    \]
    cuya inversa es $\exp:(\mathbb{R},+)\to(\mathbb{R}^+,\cdot)$.
\end{enumerate}

\medskip

La propiedad definitoria es:
\[
\log(ab)=\log(a)+\log(b),
\qquad
\log(1)=0.
\]

La multiplicación corresponde a la adición; el elemento neutro multiplicativo
corresponde al elemento neutro aditivo; los inversos multiplicativos corresponden
a los inversos aditivos.
Ninguna relación algebraica se pierde --- solo cambia su apariencia.

\medskip

No se trata de un artificio numérico ni de un atajo computacional.
Es una equivalencia estructural: los dos grupos son, desde el punto de vista
algebraico, el mismo objeto con distinta notación.

% ------------------------------------------------------------------
%\section*{Manifestación Física: la Regla de Cálculo}
% ------------------------------------------------------------------

\medskip

Este isomorfismo cobra existencia física en la regla de cálculo.
Las distancias sobre una escala logarítmica representan valores logarítmicos.
Sumar distancias equivale a multiplicar números.

\medskip

La construcción de escalas logarítmicas consiste en aplicar la transformación
\[
x \mapsto \log_{10}(x),
\]
situando cada número a una distancia proporcional a su logaritmo.
Razones iguales corresponden a distancias iguales.
La identidad estructural $\log(ab)=\log(a)+\log(b)$ se convierte en una
operación mecánica: al deslizar una regla sobre otra, la multiplicación se
realiza mediante una adición de longitudes.

\medskip

Las longitudes no son números, pero obedecen las mismas reglas estructurales.
Al explotar esta correspondencia, la regla de cálculo se convierte en un
computador analógico --- un modelo físico de un isomorfismo algebraico.

% ------------------------------------------------------------------
%\section*{El Logaritmo como Área: Interpretación Geométrica}
% ------------------------------------------------------------------

\medskip

Los logaritmos admiten también una definición puramente geométrica,
independiente de la exponenciación.
Por construcción,
\[
\ln(a)=\int_1^a \frac{1}{x}\,dx,
\]
esto es, el área con signo bajo la curva $y=\frac{1}{x}$ entre $1$ y $a$.

\medskip

Esta definición integral hace transparente la propiedad de isomorfismo.
Para $a,b>0$,
\[
\ln(ab)
=\int_1^{ab}\frac{1}{x}\,dx
=\int_1^{a}\frac{1}{x}\,dx+\int_a^{ab}\frac{1}{x}\,dx.
\]
El cambio de variable $x=at$ en la segunda integral da
\[
\int_a^{ab}\frac{1}{x}\,dx=\int_1^{b}\frac{1}{t}\,dt=\ln(b),
\]
y por tanto $\ln(ab)=\ln(a)+\ln(b)$.
El isomorfismo de grupos es un teorema, no un supuesto.

\medskip

La definición integral también proporciona de inmediato la fórmula fundamental
de derivación:
\[
\frac{d}{da}\ln(a)=\frac{1}{a}.
\]

% ------------------------------------------------------------------
%\section*{Tasa de Cambio y Aproximación Lineal}
% ------------------------------------------------------------------

\medskip

El comportamiento de $\frac{d}{da}\ln(a) = \frac{1}{a}$ conecta el logaritmo
con el marco general de las derivadas y la linealización local.

\medskip

Recordemos que para una función $y = f(x)$, la derivada en un punto $a$ es
\[
f'(a) = \lim_{h\to0}\frac{f(a+h)-f(a)}{h},
\]
la tasa de cambio instantánea y, geométricamente, la pendiente de la recta
tangente.
En un entorno de $a$, toda función diferenciable se comporta aproximadamente
como su recta tangente:
\[
f(x) \approx f(a) + f'(a)(x-a).
\]

\medskip

Esta linealidad local es el reflejo, en el ámbito del análisis, del paso de
los tensores a sus componentes tratado en el Capítulo~II: la estructura global
no lineal se vuelve manejable al aproximarse por su mejor modelo lineal en
cada punto.
La derivada sustituye localmente una aplicación no lineal por una lineal;
el logaritmo sustituye globalmente una estructura multiplicativa por una
aditiva.
Ambos son instancias del mismo principio organizador.

% ------------------------------------------------------------------
%\section*{Cálculo Numérico: Newton--Raphson Aplicado al Logaritmo}
% ------------------------------------------------------------------

\medskip

La derivada del logaritmo también permite su cálculo numérico eficiente
mediante el método de Newton--Raphson.
Calcular $\ln(x)$ equivale a resolver la ecuación
\[
e^y = x,
\]
o, de forma equivalente, encontrar el cero de $f(y) = e^y - x$.

\medskip

La iteración de Newton--Raphson reemplaza la ecuación no lineal por una
sucesión de ecuaciones lineales.
Partiendo de una aproximación inicial $y_n$, se linealiza $f$ en torno a $y_n$
y se resuelve:
\[
0 = f(y_n) + f'(y_n)(y_{n+1}-y_n).
\]
Como $f'(y) = e^y$, se obtiene
\[
y_{n+1} = y_n - \frac{e^{y_n}-x}{e^{y_n}}
         = y_n + \frac{x - e^{y_n}}{e^{y_n}}.
\]

Partiendo de $y_0 = 1{,}5$, las iteraciones sucesivas convergen rápidamente a
$\ln(5) \approx 1{,}60944$.

\medskip

La simplicidad de esta iteración no es casual.
Trabajar en base $e$ garantiza que $\frac{d}{dy}e^y = e^y$,
lo que minimiza la carga algebraica y mejora la estabilidad numérica.
Los logaritmos en otras bases se obtienen mediante la fórmula de cambio de base:
\[
\log_{b}(x) = \frac{\ln(x)}{\ln(b)}.
\]

% ------------------------------------------------------------------
%\section*{Algoritmo Histórico: Briggs y la Acumulación Binaria}
% ------------------------------------------------------------------

\medskip

La misma idea estructural subyace a los algoritmos históricos de construcción
de tablas de logaritmos.
Henry Briggs calculó tablas en base~10 acumulando pequeñas correcciones
multiplicativas; un método estrechamente relacionado, adaptado a la aritmética
binaria, se describe en~\cite{knuth}.

\medskip

El algoritmo construye el número objetivo mediante multiplicaciones sucesivas
de la forma
\[
\frac{10^k}{10^k-1},
\]
acumulando los logaritmos correspondientes cada vez que el producto parcial no
supera el objetivo.
Lo que aparece como cómputo es, en rigor, preservación de estructura:
cada paso multiplicativo corresponde a un incremento aditivo en el lado
logarítmico, ejecutando fielmente el isomorfismo $\log(ab) = \log(a) + \log(b)$.

% ------------------------------------------------------------------
%\section*{La Constante $e$ y sus Representaciones}
% ------------------------------------------------------------------

\medskip

La base natural $e$ no es una elección arbitraria.
Es la única base para la que $\frac{d}{dx}e^x = e^x$,
y la única solución a $\int_1^e \frac{1}{x}\,dx = 1$.
Su carácter aritmético queda revelado por varias construcciones independientes.

\medskip

Una definición clásica proviene del interés compuesto:
\[
e = \lim_{n\to\infty}\left(1+\frac{1}{n}\right)^n.
\]

Otra la proporciona su desarrollo en fracción continua:
\[
e = [2;\,1,2,1,1,4,1,1,6,\dots],
\]
que en forma extendida se escribe
\[
e
=
2 + \cfrac{1}{1 + \cfrac{1}{2 + \cfrac{1}{1 + \cfrac{1}{1 + \cfrac{1}{4 + \cdots}}}}}.
\]

Estas representaciones convergen rápidamente y resultan adecuadas para el
cómputo, pero su significado más profundo es estructural:
$e$ es la unidad natural del isomorfismo entre $(\mathbb{R}^+,\cdot)$
y $(\mathbb{R},+)$.

% ------------------------------------------------------------------
%\section*{Dos Linealizaciones Complementarias}
% ------------------------------------------------------------------

\medskip

El logaritmo transforma el crecimiento no lineal en desplazamiento lineal.
Convierte la multiplicación en adición, los productos en sumas
y las escalas exponenciales en escalas lineales.

\medskip

Se hacen así visibles dos manifestaciones complementarias del principio
invariante.
El logaritmo es un isomorfismo \emph{global}:
establece una equivalencia estructural perfecta entre dos mundos algebraicos,
válida en todo $\mathbb{R}^+$.
La derivada, en cambio, proporciona una linealización \emph{local}:
aproxima una transformación arbitraria por su mejor modelo lineal en cada punto.

\medskip

En ambos casos, la complejidad se hace inteligible al expresarse mediante
relaciones lineales.
Por eso los logaritmos aparecen simultáneamente en el análisis, la geometría,
los métodos numéricos y los instrumentos de cómputo.
No pertenecen a un único dominio.
Preservan estructura a través de todos ellos.